\documentclass{situationist-preprint}

\title{Curiosity as an Epistemic Virtue}
\msid{EPI-22-a}
\msfamily{EPI-22}
\msdiscipline{Epistemology}
\mstype{Principal manuscript}
\mskeywords{curiosity; intellectual virtue; inquiry; significance; trivial truths; attention; curiositas}

\begin{document}
\maketitle

\begin{abstract}
Curiosity is praised as the engine of inquiry and was condemned, in the tradition from Augustine to Aquinas, as a vice. We argue that both verdicts are right of different dispositions. Curiosity is an epistemic virtue when it is directed at significant questions, and indiscriminate inquiry is not virtuous. The argument uses a motivational account of intellectual virtue, on which a virtue involves concern for epistemic goods proportioned to their value, together with two premises: that the epistemic value of an answer depends on the significance of the question it answers, and that inquiry consumes attention, which is finite. Curiosity that tracks significance satisfies the motivational condition and reliably yields valuable knowledge. Curiosity that is indifferent to significance fails the motivational condition, whatever truths it produces, and incurs a cost in significant knowledge foregone. We show that the account accommodates the unpredictability of significance, since calibrated exploration is itself directed at expected significance, and that it does not make significance merely a matter of what the inquirer happens to care about.
\end{abstract}

\section{The problem}

\citet[X.35]{Augustine} counted curiosity, the \enquote{lust of the eyes}, among the temptations, since it seeks to know for the sake of knowing what does not concern us. \citet[II-II, q.~166--167]{Aquinas} distinguished the virtue of studiousness, which moderates the desire for knowledge, from the vice of curiosity, which directs it at unprofitable objects or pursues it in disordered ways. \citet[2.3.10]{Hume} treated curiosity, or the love of truth, as a passion whose satisfaction requires that the truth discovered be of some importance, comparing it to the pleasure of the hunt, which depends on the utility of the quarry as well as on the exercise of skill. Contemporary virtue epistemologists list curiosity or inquisitiveness among the intellectual virtues \citep{Zagzebski1996,Baehr2011}, and \citet{Inan2012} has given an account of curiosity as awareness of one's own ignorance directed at a particular object. \citet{Whitcomb2010} argued that curiosity is a desire for knowledge, and that its satisfaction conditions make knowledge the goal of inquiry.

The question is whether curiosity as such is a virtue. An inquirer who counts the blades of grass in a park, or memorizes the telephone directory, is curious in a perfectly good sense: he wants to know, and he inquires to find out. Rawls's grass counter, a man whose only pleasure is to count blades of grass in park squares and on lawns, was offered as an example of a rational plan of life \citep[\S 65]{Rawls1971}; but no one has offered him as an example of intellectual virtue. We argue that curiosity is a virtue when, and only when, it is directed at significant questions.

\section{Framework}

\begin{definition}[Curiosity]\label{d:curiosity}
\emph{Curiosity} is a disposition to desire, for their own sake, answers to questions, and to inquire in pursuit of them.
\end{definition}

\begin{definition}[Intellectual virtue]\label{d:virtue}
A disposition is an \emph{intellectual virtue} if (i) its exercise manifests a concern for epistemic goods, such as knowledge and understanding, that is proportioned to their value, and (ii) its exercise is reliably conducive to such goods in the circumstances in which it operates.
\end{definition}

Definition~\ref{d:virtue} follows the motivational account of \citet{Zagzebski1996}, on which a virtue has a motivational component and a component of reliable success, with one specification: the motivation must be proportioned to the value of its object. This is the Aristotelian thought that a virtue is a disposition to be moved by the right objects to the right degree, applied to the epistemic case.

\begin{definition}[Significance]\label{d:significance}
A question is \emph{significant} to the degree that its answer contributes to understanding, by connecting with and explaining answers to other questions, or bears on deliberation that matters to the inquirer or to her community.
\end{definition}

Definition~\ref{d:significance} is close to the notion of significance that \citet{Kitcher2001} placed at the centre of the philosophy of science. It is graded, and partly relative to the inquirer, since what bears on her deliberation depends on her situation; but its first clause is not relative to anyone's interests, since whether an answer explains others is a matter of the structure of the subject.

\begin{principle}[Value and significance]\label{p:value}
The epistemic value of knowing the answer to a question increases with the significance of the question; knowing the answer to an insignificant question has little epistemic value.
\end{principle}

Principle~\ref{p:value} is the lesson of the problem of trivial truths. If every true belief had equal epistemic value, one would do as well, epistemically, to memorize the directory as to learn how vaccines work; and since the directory contains more truths, better. \citet{Treanor2013} argued that this shows that the value of a belief is not measured by the number of truths it contains but by how much of the world it captures, which comes to much the same thing as Definition~\ref{d:significance}.

\begin{assumption}[Finite attention]\label{a:finite}
Inquiry consumes attention and time, which are finite; to inquire into one question is to forgo inquiry into others.
\end{assumption}

By EPI-11, attention is under voluntary control, while belief is not. The allocation of attention is therefore the point at which the curious inquirer exercises agency, and the point at which her curiosity can be assessed.

\section{The argument}

\begin{theorem}[Directed curiosity is a virtue]\label{t:virtue}
Curiosity whose strength, for each question, increases with the significance of the question is an intellectual virtue.
\end{theorem}

\begin{proof}
Condition (i) of Definition~\ref{d:virtue}: such curiosity is a desire for knowledge of answers, which are epistemic goods, and its strength increases with significance, which by Principle~\ref{p:value} is what their value increases with; its concern is therefore proportioned to value. Condition (ii): by Assumption~\ref{a:finite}, attention allocated to one question is withdrawn from others, and a disposition that allocates more attention to more significant questions yields, in any circumstances in which inquiry tends to succeed, more knowledge of significant answers than one that allocates it otherwise. It is therefore reliably conducive to epistemic goods weighted by their value.
\end{proof}

\begin{theorem}[Indiscriminate inquiry is not virtuous]\label{t:indisc}
Curiosity whose strength does not depend on the significance of the question is not an intellectual virtue.
\end{theorem}

\begin{proof}
Such curiosity is moved equally by questions whose answers have great and little epistemic value, and so its concern is not proportioned to value; condition (i) of Definition~\ref{d:virtue} fails. It may produce many true beliefs, but by Assumption~\ref{a:finite} each act of attention it directs to an insignificant question is withdrawn from a significant one, and so, compared with the directed disposition of Theorem~\ref{t:virtue}, it yields less of what is epistemically valuable.
\end{proof}

The second theorem reconstructs Aquinas's distinction. Studiousness, on his account, is not a weaker desire for knowledge than curiosity but a better ordered one, and the vice he called curiosity is a desire for knowledge of things that do not repay it. It also explains why the grass counter is not a model of intellectual virtue even if his plan is rational. Rawls's point concerned what a rational person might choose to do with his life; ours concerns what kind of concern for knowledge deserves admiration. The counter's inquiry may harm no one, and his beliefs about the number of blades may be true and well founded; but he is moved as much by a question whose answer explains nothing and serves no deliberation as he would be by any other, and the disposition that moves him is not proportioned to the value of what it pursues.

\section{Objections}

\begin{objection}[Unpredictable significance]
The significance of a question often becomes apparent only after it has been answered. Fleming's attention to a contaminated culture plate was not directed at a question known to be significant. A curiosity confined to recognized significance would have missed the most important discoveries.
\end{objection}

\begin{reply}
What a disposition can track is expected significance, given what the inquirer knows. A question whose answer might connect with many others, such as why bacteria failed to grow near a mould, has high expected significance even before the connection is established, and a disposition to explore anomalies is a disposition directed at such questions. Calibrated exploration of this kind, which allocates some attention to questions of uncertain significance in proportion to the chance that they will prove significant, satisfies Definition~\ref{d:virtue}. What the account excludes is not exploration but indifference: attention that would be allocated in the same way whatever the expected significance of the question.
\end{reply}

\begin{objection}[Subjectivity]
Significance is relative to interests. To the grass counter, the number of blades is significant. The account then makes any curiosity virtuous for the person who has it.
\end{objection}

\begin{reply}
Definition~\ref{d:significance} has two clauses, and neither is satisfied by the counter's question merely because he wants the answer. The number of blades explains nothing else, and it bears on no deliberation of his or anyone's, since nothing he does depends on it. Interest that is not grounded in either clause does not confer significance; it is simply the curiosity whose virtue is in question.
\end{reply}

\begin{objection}[Curiosity is a passion]
Curiosity is a desire, and a desire is not a virtue. At most it is the material on which a virtue such as studiousness operates.
\end{objection}

\begin{reply}
On the motivational account, a virtue includes a disposition of desire as one of its components. One may reserve \enquote{curiosity} for the raw desire and call its well-ordered form by another name, as Aquinas did; the substance of the thesis is unaffected. The disposition we have called directed curiosity is a virtue; the disposition to desire answers indiscriminately is not.
\end{reply}

\begin{objection}[No harm in trivia]
Where attention is not scarce, as for someone with time to spare, indiscriminate inquiry costs nothing and produces truths. Why deny it the status of a virtue?
\end{objection}

\begin{reply}
Theorem~\ref{t:indisc} does not rest on the cost alone. Even where nothing is foregone, condition (i) fails: a concern for knowledge that is equally moved by the trivial and the important is not a concern proportioned to value, and it is the proportion, not the quantity, of concern that makes a disposition admirable. Indiscriminate inquiry may be harmless, and may even be pleasant; it is not thereby virtuous.
\end{reply}

\section{Conclusion}

Curiosity is an epistemic virtue when it is directed at significant questions, because concern for knowledge proportioned to the significance of what is known is concern proportioned to its value, and such concern, given finite attention, reliably yields valuable knowledge. Indiscriminate inquiry, moved equally by every question, fails the motivational condition of virtue whatever truths it produces. The medieval distinction between studiousness and curiosity marks the same line.

\begin{thebibliography}{99}\small
\bibitem[Aquinas(1274)]{Aquinas} Thomas Aquinas (1265--1274). \emph{Summa Theologiae}.
\bibitem[Augustine(397)]{Augustine} Augustine (397--400). \emph{Confessions}.
\bibitem[Baehr(2011)]{Baehr2011} Baehr, J. (2011). \emph{The Inquiring Mind: On Intellectual Virtues and Virtue Epistemology}. Oxford University Press.
\bibitem[Hume(1739)]{Hume} Hume, D. (1739--1740). \emph{A Treatise of Human Nature}.
\bibitem[Inan(2012)]{Inan2012} Inan, I. (2012). \emph{The Philosophy of Curiosity}. Routledge.
\bibitem[Kitcher(2001)]{Kitcher2001} Kitcher, P. (2001). \emph{Science, Truth, and Democracy}. Oxford University Press.
\bibitem[Rawls(1971)]{Rawls1971} Rawls, J. (1971). \emph{A Theory of Justice}. Harvard University Press.
\bibitem[Treanor(2013)]{Treanor2013} Treanor, N. (2013). The measure of knowledge. \emph{No\^us}, 47(3), 577--601.
\bibitem[Whitcomb(2010)]{Whitcomb2010} Whitcomb, D. (2010). Curiosity was framed. \emph{Philosophy and Phenomenological Research}, 81(3), 664--687.
\bibitem[Zagzebski(1996)]{Zagzebski1996} Zagzebski, L.~T. (1996). \emph{Virtues of the Mind: An Inquiry into the Nature of Virtue and the Standard Ethical Concern of Knowledge}. Cambridge University Press.
\end{thebibliography}

\end{document}
